\section{Proof of the global Hilbert--Schmidt comparison}
\label{app:comparison}

We prove Lemma~\ref{lem:hs-global}, retaining a positive regularizer until the last step.
All traces on finite matrix algebras and their corners are unnormalized.
Set $K:=\id+\Psi$ and
\[
 \ell_\varepsilon(a):=\tfrac12\varphi_m
       \log((a+T_\Psi)^2+\varepsilon^21_m),\qquad
 F_\varepsilon(a):=\ell_\varepsilon(a)
       -\tfrac1{2m}\Tr(aK^{-1}(a)),\quad a=a^*.
\]
We will establish the uniform bound
\begin{equation}\label{eq:hs-regularized-global}
 F_\varepsilon(a)\le-\tfrac12+
              \frac{\varepsilon+\varepsilon^2/2}{1-\theta},
 \qquad \theta=\|\Psi\|_{2\to2}<1.
\end{equation}

\subsection{An inverse comparison}
If $A=A^*$ has $\|A\|\le\theta<1$ on a finite-dimensional Hilbert space $\mathcal H$ and $J=J^*=J^{-1}$, then
\begin{equation}\label{eq:hs-involution-inverse}
 L:=(J-A)^{-1}\le(\id-A)^{-1},\qquad
 \|L\|\le(1-\theta)^{-1}.
\end{equation}
The norm bound follows by factoring $J-A=J(\id-JA)$ and using the Neumann series.
To prove the order bound, put $D_0:=\id-A>0$ and $P:=(\id-J)/2$, an orthogonal projection.
When $P\ne0$, block elimination, or direct multiplication, gives
\[
 (D_0-2P)^{-1}=D_0^{-1}+2D_0^{-1}P
   (\id_{P\mathcal H}-2PD_0^{-1}P)^{-1}PD_0^{-1}.
\]
Because $D_0^{-1}\ge(1+\theta)^{-1}\id>\tfrac12\id$, the middle inverse is strictly negative on $P\mathcal H$.
The correction is negative semidefinite, which proves \eqref{eq:hs-involution-inverse}.
For $P=0$ the correction is absent.

There is an additional positivity-preservation fact when $A$ is a completely positive map and $J(x)=uxu$ with $u=u^*=u^{-1}$:
\begin{equation}\label{eq:hs-inverse-cp}
 L=\sum_{k\ge0}(JA)^kJ
\end{equation}
is completely positive.
Each summand is completely positive and the series converges in norm.
This does not assert that $L$ is positive as a Hilbert--Schmidt operator; the order comparison in \eqref{eq:hs-involution-inverse} is a separate assertion.

\subsection{A maximum over the constant term}
For fixed $\varepsilon>0$, the function $F_\varepsilon$ attains its global maximum on the real vector space of Hermitian matrices.
Indeed, $K^{-1}\ge(1+\theta)^{-1}\id$, whereas
\[
 \ell_\varepsilon(a)\le\tfrac12
       \log((\|a\|+\|T_\Psi\|)^2+\varepsilon^2).
\]
The negative quadratic term consequently dominates as $\|a\|_2\to\infty$.
Continuity gives a maximum.

Fix a maximizing $a$, and write
\[
 R:=E_m((a+T_\Psi-i\varepsilon1_m)^{-1})=H+iQ,
 \qquad H=H^*,\quad Q>0.
\]
Differentiation under the trace gives $d\ell_\varepsilon(a)[h]=m^{-1}\Tr(Hh)$ for $h=h^*$.
For example, differentiate $(a+T_\Psi)^2+\varepsilon^21_m$ and use the traced logarithmic derivative and cyclicity.
Stationarity therefore gives $H=K^{-1}(a)$.
The Dyson equation becomes
\begin{equation}\label{eq:hs-maximal-dyson}
 R^{-1}=a-\Psi(R)-i\varepsilon1_m=H-iB,
 \qquad B:=\Psi(Q)+\varepsilon1_m>0.
\end{equation}
The logarithmic primitive \eqref{eq:selfadjoint-potential} gives
\[
 \ell_\varepsilon(a)=-\frac1m\log|\det R|
       +\frac1{2m}\Tr(H\Psi(H)-Q\Psi(Q)).
\]
Taking real traces in $RR^{-1}=1_m$ yields $\Tr H^2+\Tr(Q\Psi(Q))+\varepsilon\Tr Q=m$.
Thus
\begin{equation}\label{eq:hs-maximal-value}
 F_\varepsilon(a)=-\tfrac12-\frac1m\log|\det R|
                         +\frac{\varepsilon}{2m}\Tr Q.
\end{equation}
We next control the last two terms, including when $H$ has a kernel.

\subsection{The support and kernel of the real part}
Multiplying \eqref{eq:hs-maximal-dyson} in both orders gives $[Q,B]+i[Q+B,H]=0$.
The first summand is skewadjoint and the second is selfadjoint, so each vanishes.
The matrix $QH-HB$ is then selfadjoint.
Separating the two parts of $H^2+QB+i(QH-HB)=1_m$ gives
\begin{equation}\label{eq:hs-support-identities}
 QB=BQ=1_m-H^2,\qquad QH=HB,\qquad HQ=BH.
\end{equation}
Let $e$ be the orthogonal projection onto the range of $H$, let $f:=1_m-e$, and put $d:=\operatorname{rank}e$, $r:=\operatorname{rank}f$.
The matrices $Q,B,R$ commute with these projections, since $Q$ and $B$ commute with $H^2$.

On the $e$-corner the matrix $H_e$ is invertible.
Set $u:=\operatorname{sign}(H_e)$, the selfadjoint unitary with eigenvalues the signs of the eigenvalues of $H_e$.
Since $Q_e$ commutes with $|H_e|$, \eqref{eq:hs-support-identities} implies $B_e=uQ_eu$.
Consequently $R_e^{-1}=uR_e^*u$, so $|\det R_e|=1$.
On the $f$-corner, $H_f=0$, and the same identities give $R_f=iQ_f$, $B_f=Q_f^{-1}$.
Therefore
\begin{equation}\label{eq:hs-kernel-determinant}
 \log|\det R|=\log\det Q_f.
\end{equation}
When a corner is absent, its determinant is one and its trace is zero.

For $v,w\in\{e,f\}$ write $\Psi_{vw}(x):=v\Psi(wxw)v$ for the maps between the corner matrix algebras $M_e:=eM_m(\C)e$ and $M_f:=fM_m(\C)f$.
On their direct sum the compressed map is selfadjoint, has norm at most $\theta$, and has mutually adjoint off-diagonal blocks.
Each block preserves complete positivity.

Assume first that $e\ne0$, and let $A:=\Psi_{ee}$, $J(x):=uxu$, $L:=(J-A)^{-1}$.
By \eqref{eq:hs-involution-inverse}--\eqref{eq:hs-inverse-cp}, $L$ is selfadjoint and completely positive, with $\|L\|\le(1-\theta)^{-1}$ and $L\le(\id-A)^{-1}$.
The $e$-part of $B=\Psi(Q)+\varepsilon1_m$ gives
\begin{equation}\label{eq:hs-support-elimination}
 Q_e=L\Psi_{ef}(Q_f)+\varepsilon L(e).
\end{equation}
On the $f$-corner define
\begin{equation}\label{eq:hs-effective-data}
 \Phi:=\Psi_{ff}+\Psi_{fe}L\Psi_{ef},\qquad
 D:=f+\Psi_{fe}(L(e)).
\end{equation}
Then $\Phi$ is completely positive and Hilbert--Schmidt selfadjoint, and $D\ge f>0$ on this corner.
Moreover $\id_f-\Phi>0$ as a Hilbert--Schmidt operator.
Indeed, the positive block operator $\id-\Psi|_{M_e\oplus M_f}$ has positive Schur complement
\[
 \id_f-\Psi_{ff}-\Psi_{fe}(\id_e-A)^{-1}\Psi_{ef}>0.
\]
Replacing $(\id_e-A)^{-1}$ by the smaller operator $L$ preserves this inequality.
In particular,
\begin{equation}\label{eq:hs-effective-trace}
 \Tr\Phi(f)<r\quad\text{if }r>0.
\end{equation}
The remaining Dyson equation and the trace of \eqref{eq:hs-support-elimination} now read
\begin{equation}\label{eq:hs-effective-equation}
 Q_f^{-1}=\Phi(Q_f)+\varepsilon D,\qquad
 \Tr Q=\Tr(DQ_f)+\varepsilon\Tr L(e).
\end{equation}
The second equality uses selfadjointness of $L$ and $\Psi_{ef}^*=\Psi_{fe}$.
If $e=0$, take $\Phi:=\Psi$, $D:=1_m$ and omit all terms containing $L$; the same conclusions hold.

\subsection{The positive-cone estimate}
Suppose $r>0$.
Lemma~\ref{lem:imaginary-minimum}, applied on the $f$-corner with covariance $\Phi$ and positive matrix parameter $\varepsilon D$, shows that $Q_f$ minimizes
\[
 -\log\det x+\varepsilon\Tr(Dx)+\tfrac12\Tr(x\Phi(x)),
 \qquad x>0\text{ on }f\C^m.
\]
Comparison with $x=f$, together with the trace of \eqref{eq:hs-effective-equation} multiplied by $Q_f$, yields
\[
 \log\det Q_f\ge\frac\varepsilon2\Tr(DQ_f)
       -\varepsilon\Tr D+\tfrac12(r-\Tr\Phi(f)).
\]
Substitute this into \eqref{eq:hs-maximal-value} and use \eqref{eq:hs-kernel-determinant}--\eqref{eq:hs-effective-equation}.
The terms containing $\Tr(DQ_f)$ cancel, and the last term is nonnegative by \eqref{eq:hs-effective-trace}.
We obtain
\begin{equation}\label{eq:hs-uniform-errors}
 F_\varepsilon(a)\le-\tfrac12+\frac\varepsilon m\Tr D
                          +\frac{\varepsilon^2}{2m}\Tr L(e).
\end{equation}
If $f=0$, \eqref{eq:hs-kernel-determinant} and \eqref{eq:hs-support-elimination} give the same bound with the $D$-term omitted.

These error terms are uniform over the maximizing matrices $a$ and the support projections of the associated matrices $H$.
Since the Hilbert--Schmidt norm of a rank-$d$ identity is $\sqrt d$, complete positivity and the norm bounds give
\[
 0\le\Tr L(e)\le\frac d{1-\theta},\qquad
 \Tr D\le r+\frac{\theta\sqrt{dr}}{1-\theta}
                  \le\frac m{1-\theta}.
\]
Thus \eqref{eq:hs-uniform-errors} implies \eqref{eq:hs-regularized-global} at every maximum and hence at every Hermitian constant term $a$.
Finally, spectral integration gives $\ell_\varepsilon(a)\downarrow\log\Delta(a+T_\Psi)$ as $\varepsilon\downarrow0$.
This proves \eqref{eq:hs-global-comparison}, including a limit of $-\infty$.
